\section{Глава~\ref{sec:acet}. Добавляем \nt{АЦЕТ}}

Три действия ацетилхолина измерены на срезах и в живом мозге, конфликт записи и чтения показан на модели книги, а то, что \nt{АЦЕТ} служит линией разрешения записи и сообщает ожидаемую неопределённость, --- гипотеза с частичной опорой на опыт.

{\footnotesize
\setlength{\tabcolsep}{3pt}\renewcommand{\arraystretch}{1.15}
\begin{longtable}{>{\raggedright}p{0.5cm}>{\raggedright}p{4.0cm}>{\raggedright}p{5.6cm}>{\raggedright}p{2.3cm}>{\raggedright\arraybackslash}p{1.7cm}}
\toprule
№ & Утверждение & Опыт и что измерено & Источник & Статус \\
\midrule\endfirsthead
\toprule
№ & Утверждение & Опыт и что измерено & Источник & Статус \\
\midrule\endhead
1 & \nt{АЦЕТ}-нейронов мозга немного, они собраны в нескольких группах, а выходы расходятся по всей коре: широковещательный канал & Окраска антителами к ферменту синтеза ацетилхолина и ретроградная метка у крысы: кора, гиппокамп, миндалина, обонятельная луковица и таламус получают ацетилхолин в основном от шести компактных групп клеток базального переднего мозга и ствола & \cite{Mesulam1983} & измерено \\
2 & Уровень ацетилхолина в коре высок при бодрствовании и низок в медленном сне & Микродиализ в коре и гиппокампе свободно движущихся кошек с записью ЭЭГ: выброс ацетилхолина в спокойном бодрствовании на $\sim100\%$, в активном --- на $\sim175\%$ выше, чем в медленном сне; в быстром сне в коре --- как при бодрствовании & \cite{Marrosu1995,Hasselmo1999} & измерено \\
3 & Смысл сигнала задаёт рецептор: у ацетилхолина есть быстрые рецепторы-каналы и медленные, запускающие реакции в клетке (утверждение~\ref{prop:receptor}) & Обзор измерений: действие ацетилхолина на возбудимость, передачу и пластичность зависит от места выброса, подтипа рецептора (никотиновые --- ионные каналы, мускариновые --- через G-белки) и клетки-получателя & \cite{Picciotto2012} & измерено \\
4 & Первое действие: ослабить внутренние связи, почти не трогая входы извне (множитель $1-a$ в~\eqref{eq:acetnet}) & Срезы обонятельной коры крысы: при $100$\,мкМ агонистов ацетилхолина потенциалы внутренних связей (слой Ib) падают более чем на $60\%$, внешнего входа (слой Ia) --- менее чем на $15\%$, в той же клетке; подавление пресинаптическое. Срезы гиппокампа (CA1): $89\%$ против $40\%$ для внутреннего и входного путей & \cite{HasselmoBower1992,HasselmoSchnell1994} & измерено \\
5 & Второе действие: усилить вход (множитель $\frac{1+a}{2}$) & Срезы неокортекса: никотиновые рецепторы усиливают синапсы входа из таламуса и не действуют на внутрикорковые; мускариновые ослабляют оба типа. Возбудимость нейронов ацетилхолин повышает (обзор) & \cite{Gil1997,Hasselmo2006} & измерено \\
6 & Третье действие: облегчить изменение весов (скорость $\eta a$) & Крыса: электрическая стимуляция базального ядра (источника ацетилхолина для коры) в паре с тоном вызывает у взрослого животного сильную перестройку карты слуховой коры под предъявляемый звук & \cite{KilgardMerzenich1998,Hasselmo2006} & измерено \\
7 & Правило Хебба для редких образов во втором уравнении~\eqref{eq:acetnet} даёт ассоциативную память, дописывающую образ по части & Расчёт ёмкости сети с разреженными образами и ковариационным правилом: при малой доле активных нейронов $p$ сеть хранит заметно больше образов, чем при $p=1/2$ & \cite{Tsodyks1988} & теория \\
8 & Запись при низком $a$ портит память: сеть записывает вспомненное, а не увиденное; при $a=0.1$ порча не слабее, чем при $a=0$ & \texttt{models/\allowbreak acet\_memory.py}: $200$ нейронов, пять образов по $20$, новый образ делит с одним из них $10$ нейронов, $10$ прогонов. При $a=0$ новый образ не запомнен, старый тянет $5.9$ чужих нейронов из $10$; при $a=0.1$ новый образ вспоминается ($0.97$), но оба сливаются: новый тянет $7.3$ чужих нейрона, старый --- $7.9$; с $a=0.2$ --- меньше одного & вывод в главе & модель \\
9 & Утверждение~\ref{prop:writeread}: нет уровня $a$, при котором одинаково хороши запись и чтение (рис.~\ref{fig:acetsim}) & Тот же скрипт, скорость обучения $\eta a$: новый образ запоминается за одно предъявление целиком при $a\ge0.9$, на $0.8$ при $a=0.75$, не запоминается при $a\le0.5$. Чтение по половине образа: $1.0$ при $a\le0.2$, $0.96$ при $a=0.3$, $0.04$ при $a=0.4$ & вывод в главе & модель \\
10 & В мозге один широковещательный провод --- \nt{АЦЕТ} --- переключает запись и чтение & Идея взята из моделей, построенных по измерениям на срезах. Самый прямой опыт --- у человека: блокатор мускариновых рецепторов скополамин ухудшает заучивание новых пар слов, сильнее всего пересекающихся со старыми, но не вспоминание уже выученных пар. Прямого измерения двух режимов сети нет & \cite{HasselmoAndersonBower1992,Atri2004} & гипотеза \\
11 & Фильтр~\eqref{eq:kalman} оптимален; вес входа и скорость обучения --- одна величина $P/R$, в установившемся режиме $P=\sqrt{qR}$ и память $\sqrt{R/q}$ (утверждение~\ref{prop:kalman}) & Опыта не требует: оптимальность фильтра Калмана--Бьюси --- классический результат теории оценивания; установившийся режим выведен в главе & \cite{KalmanBucy1961} & теория \\
12 & \nt{АЦЕТ} сообщает сети величину, подобную $P$, --- ожидаемую неопределённость & Предложено в вероятностной модели внимания. Прямого измерения того, что уровень ацетилхолина следует за неопределённостью оценки, нет & \cite{YuDayan2005} & гипотеза \\
13 & Часть \nt{АЦЕТ}-нейронов отвечает на награду и наказание за два десятка миллисекунд, тем сильнее, чем неожиданнее подкрепление & Мышь, оптогенетически опознанные холинергические нейроны базального переднего мозга: ответ на награду и наказание через $18\pm3$\,мс, величина растёт с неожиданностью подкрепления, ответы схожи у нейронов двух ядер & \cite{Hangya2015} & измерено \\
14 & \nt{НОРА} замечает неожиданную перемену мира, \nt{АЦЕТ} держит сеть в режиме записи & Разделение труда предложено в модели. Косвенно: у человека диаметр зрачка, связанный с \nt{НОРА}, отслеживает перемены и скорость обучения. Прямой проверки пары \nt{НОРА}--\nt{АЦЕТ} нет & \cite{YuDayan2005,Nassar2012} & гипотеза \\
15 & В медленном сне сеть в режиме чтения проигрывает записанное днём, и эти повторы переписывают его в долговременную память & Записи ансамблей нейронов гиппокампа крысы: клетки, работавшие вместе днём, чаще срабатывают вместе в последующем медленном сне и в том же порядке --- измерено. Для перезаписи: подавление рябей гиппокампа после обучения ухудшает пространственную память, а удлинение спонтанных рябей её улучшает; что причина --- низкий уровень \nt{АЦЕТ}, не доказано & \cite{WilsonMcNaughton1994,Skaggs1996,Girardeau2009,FernandezRuiz2019,Hasselmo1999} & гипотеза \\
\bottomrule
\end{longtable}}
